% 分野別類題: 線積分（ベクトル解析の初歩）
% パラメータ表示による計算・グリーンの定理・保存場とポテンシャルの3問。
% コンパイル: TEXINPUTS="../../../00_common:$TEXINPUTS" latexmk -xelatex problems.tex
\documentclass[a4paper,11pt]{article}
\usepackage{preamble}

\begin{document}

\kamokumei{類題演習 --- 線積分（Line Integrals）}

\shijibun{以下の各問に答えよ。ベクトル場 $\bm{F} = (P, Q)$ に対する線積分は $\displaystyle\int_{C} \bm{F}\cdot d\bm{r} = \int_{C} P\,dx + Q\,dy$ で定義される。曲線 $C$ がパラメータ表示 $x = x(t)$, $y = y(t)$ $(a \le t \le b)$ で与えられるとき、$dx = x'(t)\,dt$, $dy = y'(t)\,dt$ を代入して $t$ の定積分に帰着させて計算せよ。}

\shomon{問1.}{ベクトル場 $\bm{F}(x,y) = (y, x^{2})$ について、曲線 $C$: $x = t$, $y = t^{2}$ $(0 \le t \le 1)$ に沿う線積分
\[
\int_{C} \bm{F}\cdot d\bm{r} = \int_{C} y\,dx + x^{2}\,dy
\]
を求めよ。}

\shomon{問2.}{$C$ を原点を中心とする半径 $2$ の円周（反時計回り）とするとき、グリーンの定理
\[
\oint_{C} P\,dx + Q\,dy = \iint_{D} \left(\frac{\partial Q}{\partial x} - \frac{\partial P}{\partial y}\right) dx\,dy
\]
を用いて、次の線積分を求めよ。
\[
\oint_{C} \left(x^{2}y\right) dx + \left(xy^{2} + 2x\right) dy
\]}

\shomon{問3.}{ベクトル場
\[
\bm{F}(x,y) = \left(2xy + \cos x,\; x^{2} + 2y\right)
\]
について、以下の問に答えよ。
\begin{enumerate}
\item $\bm{F}$ が保存場（$\bm{F} = \nabla \phi$ となるポテンシャル $\phi$ が存在する場）であることを、$xy$ 平面全体で $\partial P/\partial y = \partial Q/\partial x$ が成り立つことを示すことで確認せよ。
\item ポテンシャル関数 $\phi(x,y)$（$\bm{F} = \nabla \phi$、すなわち $\phi_x = P$, $\phi_y = Q$）を求めよ。
\item 経路によらないことを用いて、点 $(0,0)$ から点 $(1,\pi)$ までの任意の滑らかな曲線 $C$ に沿う線積分 $\displaystyle\int_{C} \bm{F}\cdot d\bm{r}$ を求めよ。
\end{enumerate}}

\end{document}
